\chapter{Dataflow scheduling}
\label{ch:sched}

\section{Node lifecycle}
\begin{center}
\begin{tikzpicture}[font=\scriptsize,node distance=16mm,>=Stealth]
 \node[fnstate](e){EMPTY};
 \node[fnstate,right=of e](w){WAITING};
 \node[fnstate,right=of w](r){READY};
 \node[fnstate,right=of r](d){DISPATCHED};
 \draw[fnarrow] (e) -- node[fnlbl,above]{register, deps$>$0} (w);
 \draw[fnarrow] (e) to[bend left=25] node[fnlbl,above]{register, deps$=$0} (r);
 \draw[fnarrow] (w) -- node[fnlbl,above]{all producers done} (r);
 \draw[fnarrow] (r) -- node[fnlbl,above]{Director accepts} (d);
\end{tikzpicture}
\end{center}
\code{DISPATCHED} is terminal (\S\ref{ch:arch}): a real, honest
consequence, not an oversight --- see the roadmap (ch.~\ref{ch:roadmap})
for the deferred slot-reclamation work item.

\section{Verified graph topologies}
\begin{tabularx}{\textwidth}{L{3.4cm} Y}
\toprule
\rowh \thd{Topology} & \thd{What it proves} \\
\midrule
Shared producer, 2 consumers & One node's completion resolves the dependency count of \emph{two} different waiting nodes independently. \\
\rowa Multiple producers, 1 consumer & A node with \code{required}$>$1 only becomes \code{READY} once \emph{every} listed producer has completed, tracked across separate wake-up events. \\
2-hop transitive diamond ($A,B$ independent; $C$ dep-$A$; $D$ dep-$B$; $E$ dep-$C,D$) & Correct cascading wake-up two hops deep --- $E$ does not fire until $C$ and $D$ have \emph{themselves} genuinely completed, not merely been marked ready. \\
\rowa Mixed-depth fan-in (node depending on both a root and a 1-hop descendant) & Dependency resolution does not assume a uniform graph depth. \\
Multilayer (8 layer-1 neurons, random INT8 data, feeding 2 layer-2 neurons reading their real shared result bytes) & Real cross-node \emph{data} forwarding through real PSRAM --- layer-2's golden values are computed from the real bytes layer-1 actually wrote, not from an independent expectation. \\
\bottomrule
\end{tabularx}
All topologies above were exercised with the real, full
\code{neural\_multiprocessor.v} (real V1 PSRAM chain, real
\code{slot\_mem\_arbiter.v}) and verified bit-exact against a software
golden model.

\section{First-free dispatch}
The Neural Director's own scheduling policy is deliberately the simplest
one that is provably correct: a fixed, lowest-index priority scan over
currently-free slots. Round-robin, least-loaded, or any fairness-aware
alternative was explicitly deferred until real measured data showed
whether it mattered (\S\ref{sec:fairness}).

\section{Measured scheduling behavior}
\label{sec:fairness}
Real per-slot data (\code{N\_SLOTS}=4, a 128-neuron dense-layer
workload) shows a striking imbalance: slots~0 and~1 each deliver 1008
real tiles, while slots~2 and~3 deliver only 16 each --- despite all four
slots reporting near-100\% ``busy'' utilization. The cause is not
unfairness in isolation: once the shared PSRAM port is saturated
(ch.~\ref{ch:mem}), there is rarely a moment where the low-index slots
are simultaneously busy \emph{and} the high-index slots have nothing to
do, so the fixed low-index-first scan keeps re-selecting the same two
slots. This is a real, measured limitation of the current scheduler,
carried into ch.~\ref{ch:roadmap} as an open item rather than patched
without first measuring whether it is worth the added complexity for
real workloads.

\section{Correctness guarantees (measured, not assumed)}
Across the full final benchmark campaign (6 workloads $\times$ 4
\code{N\_SLOTS} configurations, re-verified after both memory
optimizations): \textbf{zero} lost jobs, \textbf{zero} duplicated jobs
(\code{jobs\_allocated == jobs\_completed == neurons\_completed} exactly,
every run), \textbf{zero} deadlocks, \textbf{zero} timeouts, correct
multi-hop dependency wake-up in every topology tested.
